\section{Project History}

\textbf{Project 18: Pedestrian Crossing Warning with Camera-Depth Fusion}

This project was developed in the framework of an academic initiative to explore the integration of camera-based and depth-sensor-based perception for autonomous driving scenarios, specifically focusing on pedestrian crossing warning systems.

\section{Development Team}

The project was developed by:
\begin{itemize}
    \item Alessandro Agosta (alessandro.agosta@studio.unibo.it)
    \item Gabriele Basigli (gabriele.basigli@studio.unibo.it)
    \item Lorenzo Tosi (lorenzo.tosi10@studio.unibo.it)
\end{itemize}

\section{Development Environment}
\subsection{Software}

\begin{verbatim}
Operating System: Ubuntu 22.04 LTS
Python Version: 3.7
CARLA Version: 0.9.15
\end{verbatim}

\section{Development Process}

\subsection{Version Control}

We used Git for version control with the following structure:
\begin{itemize}
    \item \texttt{main} branch: Stable releases, production-ready code
    \item \texttt{develop} branch: Integration branch for ongoing development
    \item Feature branches: \verb|feature/<name>| for individual features
    \item Issue branches: \verb|issue/<number>| for bug fixes
\end{itemize}

\subsection{Workflow}

Our development followed an agile approach with:
\begin{enumerate}
    \item Work on \texttt{develop} branch via feature branches
    \item Regular pull requests to \texttt{develop}
    \item Periodic merges to \texttt{main} for releases
    \item Issue tracking via GitHub Issues
\end{enumerate}

\subsection{Code Organization}

The project is organized into packages under \texttt{src/}: \texttt{core/} (simulation engine, perception, state), \texttt{controllers/} (vehicle control), \texttt{sensors/} (sensor lifecycle), \texttt{utils/} (collision, detection, Kalman, visualization), and \texttt{scenarios/} (scenario definitions). A detailed description of each module is provided in the Design chapter.

\subsection{Key Implementation Decisions}

\begin{enumerate}
    \item \textbf{CARLA Integration}: Used CARLA's TCP-IP interface for the simulation environment, allowing external control and sensor data access.

    \item \textbf{YOLOv8 for Detection}: Chose YOLOv8 over other detection frameworks due to its real-time performance and ease of integration with PyTorch.

    \item \textbf{Kalman Filter for Tracking}: Implemented a custom 2D Kalman filter for pedestrian tracking instead of using external tracking libraries, for better control and understanding of the algorithm.

    \item \textbf{TTC for Collision}: Used Time-to-Collision as the collision metric, which is a standard metric in the autonomous driving community.

    \item \textbf{Pygame for Visualization}: Selected Pygame for real-time visualization due to its simplicity and cross-platform compatibility.

    \item \textbf{YAML Configuration}: Used YAML for configuration due to its human-readability and ease of editing.
\end{enumerate}

\section{Testing}

\subsection{Test Structure}

Tests are organized in the \texttt{tests/} directory:

\begin{itemize}
    \item \texttt{test\_collision.py}: Core collision risk scenarios (4 tests)
    \item \texttt{test\_collision\_edge\_cases.py}: Boundary conditions for collision logic (18 tests)
    \item \texttt{test\_kalman.py}: Kalman filter initialization, convergence, and ego-motion compensation (11 tests)
    \item \texttt{test\_scenarios.py}: Scenario loading, validation, and parsing (26 tests)
    \item \texttt{test\_perception.py}: Perception pipeline with mocked detections (8 tests)
    \item \texttt{test\_state.py}: SimulationState defaults and mutability (12 tests)
    \item \texttt{test\_config.py}: YAML configuration loading (6 tests)
    \item \texttt{test\_sensor\_manager.py}: Sensor lifecycle management (8 tests)
    \item \texttt{test\_utils.py}: Logger utility (1 test)
\end{itemize}

\subsection{Test Coverage}

The test suite contains 98 tests covering all pure logic modules. Modules that directly interact with the CARLA simulator (e.g., \texttt{engine.py}, \texttt{pedestrian.py}, \texttt{carla\_utils.py}) are not unit-tested as they require a running simulator instance.

\begin{itemize}
    \item Collision module: full coverage (22 test cases including edge cases)
    \item Kalman filter: full coverage (initialization, convergence, noise smoothing)
    \item Scenario loading and validation: full coverage (30 tests)
    \item Perception system: covered via mocked detections and depth images
    \item Configuration and state: full coverage
    \item Sensor manager: full coverage via mocked CARLA actors
\end{itemize}

\subsection{Test Execution}

\begin{verbatim}
# Run all tests
pytest

# Run specific test file
pytest tests/test_collision.py

# Run with verbose output
pytest -v

# Run with coverage
pytest --cov=src --cov-report=html
\end{verbatim}

\section{Deployment}

\subsection{Requirements}

To run the project, follow these steps:

\begin{verbatim}
# Create conda environment
conda create -n carla-env python=3.11

# Activate environment
conda activate carla-env

# Install CARLA
pip install git+https://github.com/carla-simulator/carla.git

# Install project dependencies
pip install -r requirements.txt
\end{verbatim}

\textbf{Additional requirements for simulation:}

\subsection{Running Simulations}

\textbf{List available scenarios:}

\begin{verbatim}
python src/run_simulation.py --list-scenarios
\end{verbatim}

\textbf{Run a specific scenario:}

\begin{verbatim}
python src/run_simulation.py \\
  --config config/config.yaml \\
  --scenario non-colliding-pedestrian
\end{verbatim}

\textbf{Use default scenario:}

\begin{verbatim}
python src/run_simulation.py --config config/config.yaml
\end{verbatim}

\textbf{Manual Control:}

To enable manual driving using the keyboard, append the \texttt{--manual} flag:

\begin{verbatim}
python src/run_simulation.py --scenario colliding-pedestrian --manual
\end{verbatim}

Key bindings:
\begin{itemize}
    \item \textbf{W / Up Arrow}: Throttle
    \item \textbf{S / Down Arrow}: Brake
    \item \textbf{A / Left Arrow}: Steer Left
    \item \textbf{D / Right Arrow}: Steer Right
    \item \textbf{Space}: Hand-brake
\end{itemize}
\textit{Note:} The Automatic Emergency Braking (AEB) system remains active in manual mode to prevent imminent collisions.

\subsection{Configuration}

The system is configured via a YAML file (\texttt{config/config.yaml}). The full configuration structure, including all available parameters for CARLA, sensors, and scenarios, is described in the Design chapter.

\subsection{Visualization}

The visualization window appears automatically during simulation.\\ Press \texttt{ESC} or click the X button to close.

\section{Continuous Integration}

We use GitHub Actions for automated testing and CI/CD:

\begin{verbatim}
.github/workflows/
    deploy.yml         # Build and deploy on tag release
\end{verbatim}

\subsection{Release Process}

For each release:
\begin{enumerate}
    \item Make changes in feature branch
    \item Test locally and commit
    \item Create pull request to \texttt{develop}
    \item Review and merge PR
    \item Merge to \texttt{main}
    \item Release via GitHub Releases UI
    \item Update report via report generation
\end{enumerate}
